\section{Background}
\label{sec:background}
This section summarizes the Cray Programming Environment, the AFAR compiler
drop model, and the Frontier system context used in this study.

\subsection{Cray Programming Environment}
\label{sec:cray-pe}
The Cray PE exposes consistent compiler drivers (\texttt{ftn}, \texttt{cc}, and
\texttt{CC}) backed by a module hierarchy. Users load a \texttt{cpe/<ver>}
module and a corresponding \texttt{PrgEnv-*} module (e.g.,
\texttt{PrgEnv-amd/\allowbreak \prgenvversion}) which configure compiler
wrappers, MPI integration, and library metadata. The wrappers embed PE-specific
assumptions about flag handling and \texttt{PKG\_\allowbreak CONFIG\_PATH}
(including the \texttt{PE\_AMD\_\allowbreak FIXED\_\allowbreak PKGCONFIG\_PATH}
convention), and most build systems assume those semantics remain stable.
The PE module hierarchy layers target modules (e.g.,
\texttt{craype-x86-trento}, \texttt{craype-accel-amd-gfx90a}), MPI, and vendor
libraries on top of the base compiler environment. This hierarchy is designed
to keep wrapper behavior consistent even when users mix additional library
modules or swap MPI implementations, which places a premium on stable metadata
paths for `pkg-config` and compiler wrappers.

\subsection{AMD Compiler and ROCm Drops}
\label{sec:amd-drops}
AFAR drops bundle AMD compilers with a matched ROCm toolkit, and they are
released on a rapid cadence. Each drop requires validation against Cray PE
libraries and MPI, but the drops themselves do not integrate with the PE
module hierarchy. Afar-prgenv provides that integration layer while preserving
PE wrapper behavior.

At the time of this evaluation, the latest AFAR drop available on Frontier is
\afarversion, paired with \texttt{rocm/\allowbreak \afarrocm} and
\texttt{amd/\allowbreak \afarrocm}.
While AFAR drops ship their own ROCm runtime, site workflows may still load a
system ROCm module (such as `rocm/\afarrocm`) to satisfy dependencies of
selected vendor libraries that lag behind the drop.
AFAR drops also package MPI module files under versioned include directories
(\texttt{include/mpich3.4a2} or \texttt{include/mpich4.3.1}) to match Cray MPICH
8.x or 9.x layouts. Maintaining compatibility with these layouts is critical
for Cray PE builds that rely on \texttt{pkg-config} to discover MPI module
paths.

\subsection{Metadata and pkg-config}
\label{sec:pkgconfig-bg}
Cray PE libraries expose their compile and link flags through `.pc` files
consumed by `pkg-config`. Build systems rely on this metadata to resolve
library search paths and include directories. Afar-prgenv relies on
`pkg-config` to keep MPI and vendor library metadata aligned with AFAR
compilers, and it uses shim `.pc` files when PE modules do not provide
compatible metadata.

\subsection{Frontier as a Target System}
\label{sec:frontier}
Frontier is an HPE Cray EX system with AMD GPUs. The system uses a
module-driven PE; at the time of this work, \texttt{cpe/\allowbreak \peversion}
and \texttt{PrgEnv-amd/\allowbreak \prgenvversion} are available. The
recommended baseline for AFAR testing in this workspace loads
\texttt{rocm/\allowbreak \afarrocm} and \texttt{amd/\allowbreak \afarrocm}
alongside \texttt{PrgEnv-amd}, then adds the AFAR module tree from
\texttt{/sw/crusher/ums/compilers/\allowbreak modulefiles/}.
The target accelerator for offload testing is AMD GFX90a. We select it via the
\texttt{craype-accel-}\allowbreak\texttt{amd-gfx90a} module in the PE stack.
